\subsection{Matrix Representation and Vector Spaces}

When we studied scalar differential equations in the previous section, we tracked how a single quantity like temperature or position evolves over time. Real physical systems, however, rarely operate in isolation. A satellite tumbling through space has six degrees of freedom---three for translation and three for rotation---all coupled together through gravitational and aerodynamic forces. An electrical power grid must balance voltage and current at thousands of interconnected nodes simultaneously. A chemical reactor must track concentrations of multiple species that react and diffuse together. In each of these cases, understanding the system requires tracking not one quantity but many, and these quantities influence each other in complex ways.

Matrices provide the mathematical language for describing such interconnected, multi-variable systems with remarkable economy and clarity. A matrix is simply a rectangular array of numbers arranged in rows and columns, but this simple structure carries enormous expressive power. When we organize the variables of a system into a vector---an ordered list of quantities---and collect the relationships between these variables into a matrix, we transform a potentially unwieldy system of equations into a single compact statement. This transformation is not merely cosmetic; it reveals underlying structure, enables powerful analytical techniques, and forms the foundation for modern control theory.

Consider a simple mechanical system to understand how matrices emerge naturally from physical principles. Imagine two masses connected by springs. The left mass $m_1$ is attached to a fixed wall by a spring with stiffness $k_1$, while the two masses are connected by a second spring with stiffness $k_2$. Let $x_1(t)$ and $x_2(t)$ denote the displacements of the two masses from their equilibrium positions, measured to the right. Applying Newton's second law to each mass yields the equations of motion.

For the first mass $m_1$, the force from the left spring is $-k_1 x_1$ (opposing the displacement), while the middle spring exerts a force $k_2(x_2 - x_1)$ that pulls $m_1$ to the right when $x_2 > x_1$. The net force on $m_1$ is therefore $-k_1 x_1 + k_2(x_2 - x_1)$, and Newton's law gives:
\[
m_1 \ddot{x}_1 = -k_1 x_1 + k_2(x_2 - x_1) = -(k_1 + k_2)x_1 + k_2 x_2.
\]

For the second mass $m_2$, the middle spring pulls to the left with force $k_2(x_1 - x_2)$, while any external force or damping would appear here. Assuming no external forces for simplicity, Newton's law yields:
\[
m_2 \ddot{x}_2 = k_2(x_1 - x_2) = k_2 x_1 - k_2 x_2.
\]

We now have two coupled second-order differential equations. To express this system in matrix form, we first introduce the acceleration vector $\ddot{\mathbf{x}} = \begin{bmatrix} \ddot{x}_1 & \ddot{x}_2 \end{bmatrix}^\top$ and the displacement vector $\mathbf{x} = \begin{bmatrix} x_1 & x_2 \end{bmatrix}^\top$. The equations can be rewritten as:
\[
m_1 \ddot{x}_1 = -(k_1 + k_2)x_1 + k_2 x_2, \quad
m_2 \ddot{x}_2 = k_2 x_1 - k_2 x_2.
\]

Defining the mass matrix $\mathbf{M} = \begin{bmatrix} m_1 & 0 \\ 0 & m_2 \end{bmatrix}$ and the stiffness matrix $\mathbf{K} = \begin{bmatrix} k_1 + k_2 & -k_2 \\ -k_2 & k_2 \end{bmatrix}$, we can write the system as a single matrix equation:
\[
\mathbf{M} \ddot{\mathbf{x}} + \mathbf{K} \mathbf{x} = \mathbf{0}.
\]

This compact representation captures both equations simultaneously. The matrices $\mathbf{M}$ and $\mathbf{K}$ encode all the physical parameters and coupling relationships between the two masses. Notice how the stiffness matrix $\mathbf{K}$ exhibits a important structural property: it is symmetric, meaning $\mathbf{K} = \mathbf{K}^\top$, which reflects the underlying physics of energy conservation in this conservative system.

The matrix-vector multiplication that underlies this representation deserves careful examination. When we multiply a matrix $\mathbf{A} \in \mathbb{R}^{m \times n}$ by a vector $\mathbf{x} \in \mathbb{R}^n$, the result is a new vector $\mathbf{y} = \mathbf{A}\mathbf{x} \in \mathbb{R}^m$ whose $i$-th component is given by the dot product of the $i$-th row of $\mathbf{A}$ with the vector $\mathbf{x}$:
\[
y_i = \sum_{j=1}^n a_{ij} x_j.
\]

Each row of the matrix represents one equation in the system, and each column represents how one input variable contributes to all the output equations. In our mechanical example, the $(1,2)$ entry of the stiffness matrix ($k_2$) tells us how the displacement of the second mass affects the force on the first mass. The symmetry of this entry with the $(2,1)$ entry reflects the fact that the spring connecting the two masses transmits force bidirectionally with equal magnitude.

To build intuition about matrix-vector multiplication, consider a simple $2 \times 2$ system that we can work out completely by hand. Suppose we have the system:
\[
\begin{aligned}
3x_1 + x_2 &= 7, \\
x_1 - 2x_2 &= -4.
\end{aligned}
\]

In matrix form, this becomes $\mathbf{A}\mathbf{x} = \mathbf{b}$ where $\mathbf{A} = \begin{bmatrix} 3 & 1 \\ 1 & -2 \end{bmatrix}$, $\mathbf{x} = \begin{bmatrix} x_1 \\ x_2 \end{bmatrix}$, and $\mathbf{b} = \begin{bmatrix} 7 \\ -4 \end{bmatrix}$. Computing the product $\mathbf{A}\mathbf{x}$ explicitly:
\[
\mathbf{A}\mathbf{x} = \begin{bmatrix} 3 & 1 \\ 1 & -2 \end{bmatrix} \begin{bmatrix} x_1 \\ x_2 \end{bmatrix} = \begin{bmatrix} 3x_1 + x_2 \\ x_1 - 2x_2 \end{bmatrix},
\]
and setting this equal to $\mathbf{b}$ recovers our original system. To solve this system, we can use Gaussian elimination or, for $2 \times 2$ systems, the explicit formula:
\[
\mathbf{x} = \mathbf{A}^{-1}\mathbf{b} = \frac{1}{\det(\mathbf{A})} \begin{bmatrix} -2 & -1 \\ -1 & 3 \end{bmatrix} \begin{bmatrix} 7 \\ -4 \end{bmatrix},
\]
where $\det(\mathbf{A}) = 3(-2) - (1)(1) = -6 - 1 = -7$ is the determinant. Computing the solution:
\[
\mathbf{x} = -\frac{1}{7} \begin{bmatrix} -2(7) + (-1)(-4) \\ -1(7) + 3(-4) \end{bmatrix} = -\frac{1}{7} \begin{bmatrix} -14 + 4 \\ -7 - 12 \end{bmatrix} = -\frac{1}{7} \begin{bmatrix} -10 \\ -19 \end{bmatrix} = \begin{bmatrix} 10/7 \\ 19/7 \end{bmatrix}.
\]
Thus $x_1 = 10/7$ and $x_2 = 19/7$. Verification shows that $3(10/7) + 19/7 = 30/7 + 19/7 = 49/7 = 7$ and $10/7 - 2(19/7) = 10/7 - 38/7 = -28/7 = -4$, confirming the solution.

The geometric interpretation of vectors and matrices provides crucial intuition that extends far beyond these algebraic calculations. A vector in $\mathbb{R}^2$ can be visualized as an arrow from the origin to a point $(x_1, x_2)$ in the plane. The length of this arrow, or the vector's norm, is given by $\|\mathbf{x}\| = \sqrt{x_1^2 + x_2^2}$. The direction of the arrow tells us the relative proportions of the two components. When we multiply a matrix by a vector, we can think of the matrix as a linear transformation that stretches, rotates, or shears the vector in characteristic ways.

Consider the matrix $\mathbf{A} = \begin{bmatrix} 2 & 0 \\ 0 & 3 \end{bmatrix}$, which scales the first coordinate by 2 and the second by 3. The vector $\mathbf{x} = \begin{bmatrix} 1 \\ 1 \end{bmatrix}$ is transformed to $\mathbf{A}\mathbf{x} = \begin{bmatrix} 2 \\ 3 \end{bmatrix}$, which points in roughly the same direction but is stretched and tilted. The matrix $\mathbf{R} = \begin{bmatrix} \cos\theta & -\sin\theta \\ \sin\theta & \cos\theta \end{bmatrix}$ rotates vectors counterclockwise by angle $\theta$. The matrix $\mathbf{S} = \begin{bmatrix} 1 & 1 \\ 0 & 1 \end{bmatrix}$ shears the plane, moving points parallel to the $x_1$-axis by an amount proportional to their $x_2$-coordinate. Each of these transformations has a distinct geometric character that we can visualize, and understanding these geometric effects prepares us to interpret more complex matrices that arise in control systems.

The concept of linear independence is fundamental to understanding vector spaces and, ultimately, to determining whether a system of equations has a unique solution. Two vectors $\mathbf{v}_1$ and $\mathbf{v}_2$ in $\mathbb{R}^2$ are linearly dependent if one is a scalar multiple of the other, meaning they point in the same or opposite directions along the same line through the origin. If no such scalar exists---if neither vector can be expressed as a multiple of the other---then they are linearly independent. Geometrically, two linearly independent vectors in $\mathbb{R}^2$ span the entire plane, meaning any vector in the plane can be written as a unique linear combination $c_1\mathbf{v}_1 + c_2\mathbf{v}_2$ for some coefficients $c_1$ and $c_2$.

To make this concrete, consider the vectors $\mathbf{v}_1 = \begin{bmatrix} 1 \\ 0 \end{bmatrix}$ and $\mathbf{v}_2 = \begin{bmatrix} 1 \\ 1 \end{bmatrix}$. These are linearly independent because there is no scalar $\alpha$ such that $\mathbf{v}_1 = \alpha\mathbf{v}_2$ (the second components would require $0 = \alpha$, but then the first components would give $1 = 0$, a contradiction). Any vector $\mathbf{x} = \begin{bmatrix} x_1 \\ x_2 \end{bmatrix}$ can be expressed as:
\[
\mathbf{x} = (x_1 - x_2)\mathbf{v}_1 + x_2\mathbf{v}_2 = (x_1 - x_2)\begin{bmatrix} 1 \\ 0 \end{bmatrix} + x_2\begin{bmatrix} 1 \\ 1 \end{bmatrix},
\]
and this representation is unique. The coefficients $(x_1 - x_2)$ and $x_2$ are determined uniquely by $x_1$ and $x_2$. Now consider what happens if we choose linearly dependent vectors, say $\mathbf{w}_1 = \begin{bmatrix} 1 \\ 0 \end{bmatrix}$ and $\mathbf{w}_2 = \begin{bmatrix} 2 \\ 0 \end{bmatrix}$. Both vectors lie along the $x_1$-axis, so they can only generate vectors in that one-dimensional subspace. The vector $\mathbf{x} = \begin{bmatrix} 0 \\ 1 \end{bmatrix}$ cannot be written as a linear combination of $\mathbf{w}_1$ and $\mathbf{w}_2$ because any such combination would have second component zero.

The dimension of a vector space is defined as the minimum number of linearly independent vectors needed to span the space. $\mathbb{R}^2$ has dimension 2 because any two linearly independent vectors span the entire plane. $\mathbb{R}^3$ has dimension 3, $\mathbb{R}^n$ has dimension $n$, and so on. This concept extends to more abstract spaces: the space of all polynomials of degree at most 2 has dimension 3 (spanned by $\{1, t, t^2\}$), and the space of $2 \times 2$ matrices has dimension 4 (spanned by the four matrices with a single 1 in each position). Understanding dimension is essential because it tells us how many independent parameters we need to describe a system completely.

When we form a matrix from the columns of linearly independent vectors, that matrix is said to be full rank, and its determinant is nonzero. For a square matrix, these three properties are equivalent: full rank, nonzero determinant, and invertibility (the existence of a matrix $\mathbf{A}^{-1}$ such that $\mathbf{A}^{-1}\mathbf{A} = \mathbf{A}\mathbf{A}^{-1} = \mathbf{I}$, where $\mathbf{I}$ is the identity matrix). If a matrix is not full rank, it is singular, its determinant is zero, and the system $\mathbf{A}\mathbf{x} = \mathbf{b}$ may have either no solution or infinitely many solutions depending on whether $\mathbf{b}$ lies in the column space of $\mathbf{A}$.

This brings us to the connection with state-space representations in control theory. In the state-space approach, we describe a dynamical system not by a single high-order differential equation but by a set of first-order differential equations. For a linear time-invariant system, these equations take the form:
\[
\dot{\mathbf{x}}(t) = \mathbf{A}\mathbf{x}(t) + \mathbf{B}\mathbf{u}(t), \quad
\mathbf{y}(t) = \mathbf{C}\mathbf{x}(t) + \mathbf{D}\mathbf{u}(t),
\]
where $\mathbf{x}(t) \in \mathbb{R}^n$ is the state vector, $\mathbf{u}(t) \in \mathbb{R}^m$ is the input vector, $\mathbf{y}(t) \in \mathbb{R}^p$ is the output vector, and $\mathbf{A}, \mathbf{B}, \mathbf{C}, \mathbf{D}$ are matrices of appropriate dimensions. The matrix $\mathbf{A}$ is called the system matrix, and its eigenvalues determine the stability and natural response characteristics of the system. The matrix $\mathbf{B}$ describes how inputs affect the state derivatives, $\mathbf{C}$ describes how states are measured as outputs, and $\mathbf{D}$ captures any direct feedthrough from inputs to outputs.

Returning to our two-mass-spring example, we can convert the second-order system into state-space form. Let us define the state vector to include both positions and velocities:
\[
\mathbf{x} = \begin{bmatrix} x_1 \\ x_2 \\ \dot{x}_1 \\ \dot{x}_2 \end{bmatrix}.
\]

Then the original equations $m_1\ddot{x}_1 = -(k_1+k_2)x_1 + k_2 x_2$ and $m_2\ddot{x}_2 = k_2 x_1 - k_2 x_2$ become four first-order equations:
\[
\dot{x}_1 = \dot{x}_1, \quad
\dot{x}_2 = \dot{x}_2, \quad
m_1\dot{\dot{x}}_1 = -(k_1+k_2)x_1 + k_2 x_2, \quad
m_2\dot{\dot{x}}_2 = k_2 x_1 - k_2 x_2.
\]

In matrix form, this is:
\[
\begin{bmatrix} \dot{x}_1 \\ \dot{x}_2 \\ \ddot{x}_1 \\ \ddot{x}_2 \end{bmatrix} =
\begin{bmatrix}
0 & 0 & 1 & 0 \\
0 & 0 & 0 & 1 \\
-\frac{k_1+k_2}{m_1} & \frac{k_2}{m_1} & 0 & 0 \\
\frac{k_2}{m_2} & -\frac{k_2}{m_2} & 0 & 0
\end{bmatrix}
\begin{bmatrix} x_1 \\ x_2 \\ \dot{x}_1 \\ \dot{x}_2 \end{bmatrix}.
\]

The $4 \times 4$ matrix that multiplies the state vector is the state-space $\mathbf{A}$ matrix for this system. Its eigenvalues, found by solving $\det(s\mathbf{I} - \mathbf{A}) = 0$, determine the natural frequencies and damping characteristics of the coupled masses. This example illustrates how matrices encode the complete dynamics of a multi-variable system in a form amenable to analysis and controller design.

Physical systems across engineering disciplines naturally give rise to matrix representations. In electrical circuits, Kirchhoff's voltage and current laws lead to systems of equations that can be organized into impedance and admittance matrices. For a circuit with $n$ nodes, the node voltage equations form an $(n-1) \times (n-1)$ system where each entry encodes the conductances connecting nodes. In chemical engineering, mass balances on reactors with multiple species lead to stoichiometric matrices whose rows represent species and columns represent reactions. The null space of this matrix reveals conservation laws and dependent species. In aerospace engineering, the equations of motion for aircraft or spacecraft involve inertia matrices that couple rotational and translational dynamics when the body axes are not aligned with principal axes.

The power of matrix methods lies not only in their compactness but in the universal techniques they enable. The same operations---matrix multiplication, inversion, eigenvalue decomposition---apply regardless of whether the matrices arise from mechanical, electrical, chemical, or economic systems. This universality allows control engineers to develop general methods that transfer across application domains. When we learn to analyze the eigenvalues of a matrix to determine stability, we learn a skill applicable to power systems, chemical processes, aerospace vehicles, and economic models alike.

As we proceed through this textbook, we will rely increasingly on matrix and vector-space methods. The state-space representation introduced here will be our primary framework for analyzing and designing control systems. Understanding how matrices encode system dynamics, how vector spaces describe possible behaviors, and how linear transformations map inputs to outputs forms the mathematical foundation for everything that follows. The geometric intuition developed in this section---visualizing vectors as arrows, understanding linear independence as spanning directions, and interpreting matrices as geometric transformations---will prove invaluable as we tackle more complex systems and design sophisticated controllers.
